%============================================================ BLOQUE 1 · Por qué (10 min)
\ugrsection{01}{¿Por qué tener tu propio LLM?}{Inicio: 0:00 · 10 min\\Soberanía, control y coste}

\ugrimageframe{imagenes/s01_infraestructura.jpeg}{El LLM, nueva infraestructura básica}{Como el correo, las bases de datos o el almacenamiento · Ilustración generada con IA}

\begin{frame}{Infraestructura, no experimento}{Como antes el correo electrónico, las bases de datos o los escritorios virtuales}
	\begin{columns}[T]
		\begin{column}{0.5\textwidth}
			\begin{itemize}
				\item Casi todo software institucional llevará, o necesitará, un modelo de lenguaje.
				\item Quien tardó en tener su propio correo o almacenamiento acabó dependiendo de terceros.
				\item La pregunta no es \emph{si} usarás IA, sino \hlugr{quién custodia los datos}.
			\end{itemize}
		\end{column}
		\begin{column}{0.46\textwidth}
			\ugrfact[Operadores, no espectadores]{Entender la tecnología, desplegarla y gobernarla como cualquier otro servicio.}
		\end{column}
	\end{columns}
	\keymsg{Un servicio LLM es infraestructura: hay que operarlo, protegerlo y escalarlo.}
\end{frame}

\begin{frame}{La soberanía es una decisión institucional}{Qué cambia cuando el modelo se ejecuta dentro del perímetro}
	\ugrcompare{Servicio comercial alojado}{%
		Cada consulta sale de la institución.\\
		Las reglas las fija el proveedor.\\
		El coste crece con el uso.\\
		Acceso a los modelos más capaces.}
	{Infraestructura propia}{%
		Los datos no salen del perímetro (RGPD, ENS).\\
		Cada institución define su política.\\
		Coste $\approx$ amortización + operación.\\
		Modelos abiertos, algo por detrás.}
	\keymsg{No es solo una decisión técnica: es decidir quién custodia la información.}
\end{frame}

\begin{frame}{El caso que seguiremos: MiIA}{Un servicio institucional de IA generativa, construido y operado en casa}
	\begin{columns}[c]
		\begin{column}{0.58\textwidth}
			\centering
			\adjustbox{max width=\linewidth, max height=3.9cm}{%
				\begin{tikzpicture}[every node/.append style={font=\footnotesize}]
					\node[rbx, minimum width=3.3cm, minimum height=1.6cm] (core) at (0,0)
					{\textbf{Núcleo MiIA}\\[2pt] NGINX $\rightarrow$ LiteLLM\\ $\rightarrow$ vLLM \,(Qwen)\\ {\scriptsize Docker}};
					\node[bx] (chat) at (-2.9,1.9)  {MiIA Chat\\ {\scriptsize Open WebUI}};
					\node[bx] (mail) at (2.9,1.9)   {ThunderAI\\ {\scriptsize Thunderbird}};
					\node[bx] (code) at (-3.5,-0.1) {Continue\\ {\scriptsize VS Code}};
					\node[bx] (oc)   at (3.5,-0.1)  {OpenCode\\ {\scriptsize terminal}};
					\node[bx] (epf)  at (0,-2.2)    {EPFReader\\ {\scriptsize documentos}};
					\foreach \n in {chat,mail,code,oc,epf} \draw[bar] (\n) -- (core);
				\end{tikzpicture}}
		\end{column}
		\begin{column}{0.36\textwidth}
			\bigstat{5}{del navegador al terminal.}\\[4pt]
			\bigstat[ugrInk]{1}{un solo endpoint OpenAI.}
		\end{column}
	\end{columns}
	\keymsg{Pila 100\,\% de código abierto, en contenedores y sobre hardware propio.}
\end{frame}
